\begin{abstract}
Feed-forward multi-view reconstruction models such as VGGT and DéjàView predict scene geometry and camera parameters in a single pass, but achieving high reconstruction quality can require substantial computational resources. We investigate whether test-time optimization (TTO) can recover the performance lost when distilling large reconstruction models into smaller students. We will train progressively compressed students from teacher predictions, refine their outputs using multi-view consistency, and compare them with the original models at matched computational budgets. The evaluation will characterize the trade-off between model size, reconstruction quality, and optimization cost, identifying which initialization errors TTO can correct and the smallest student that approaches teacher-level performance.

\begin{comment}
\begin{itemize}
    \item Feed-forward multi-view transformers (VGGT, DéjàView) predict cameras, depth and point maps in a single pass, but state-of-the-art models have grown to $\sim$1B parameters
    \item Two observations: compact models can be competitive (DéjàView, 117M), and optimization on top of feed-forward predictions still gives large gains (VGGT + BA)
    \item Question: how small can a feed-forward model become if test-time optimization (TTO) refines its output? Does high-quality reconstruction need an accurate initial prediction?
    \item Approach: distill VGGT/DéjàView into progressively smaller students, refine their predictions with TTO, and compare against the teachers
    \item Expected result: quality-vs-compute Pareto curves and a characterization of which initialization errors TTO can and cannot correct
\end{itemize}
\end{comment}

\end{abstract}